\documentclass{article}

% Required for row colouring in tables, including if xcolor is loaded by the style.
\PassOptionsToPackage{table}{xcolor}

% Official NeurIPS 2026 style for double-blind workshop submission.
% Do NOT use [final] or [preprint] for the revised anonymous submission.
\usepackage[dblblindworkshop]{neurips_2026}

% Standard packages
\usepackage[utf8]{inputenc}
\usepackage[T1]{fontenc}
\usepackage{hyperref}
\usepackage{url}
\usepackage{xurl}
\usepackage{amsmath,amssymb,amsfonts}
\usepackage{graphicx}
\usepackage{booktabs}
\usepackage{xcolor}
\usepackage{nicefrac}
\usepackage{microtype}
\usepackage{enumitem}
\usepackage{multirow}
\usepackage{caption}
\usepackage{subcaption}
\usepackage{CJKutf8}

% Colours for tables
\definecolor{rowred}{RGB}{255,242,242}
\definecolor{rowgreen}{RGB}{242,255,242}

% Caption style
\captionsetup{font=small,labelfont=bf,skip=4pt}

% Figure helper
% Keep your image files in the same project folder with these exact names.
% Example: figure1_layer_suppression.png
\newcommand{\fig}[2]{%
  \includegraphics[width=\linewidth]{#1}%
}

% Paper title
\title{Last-Layer Suppression: A Consistent Causal Mechanism\\
of Factual Hallucination in Transformer Language Models}

% Required for NeurIPS workshop format
\workshoptitle{ICML 2026 Workshop on Mechanistic Interpretability}

% Anonymous double-blind author block
\author{Anonymous Author(s)\\
Affiliation withheld for double-blind review}

\begin{document}

\maketitle

\begin{abstract}
Factual hallucinations in large language models are widely attributed to absent
knowledge. We present mechanistic and causal evidence that a substantial fraction
instead reflect \emph{retrieval failures}: correct factual knowledge is encoded in
intermediate transformer layers but systematically suppressed at the final output
layer, a mechanism we term \textbf{Last-Layer Suppression} (LLS). We provide six
complementary lines of evidence spanning \textbf{12 models across 6 architecture
families (124M--8B parameters)}, validated on \textbf{680 prompts} covering 150
facts, 10 knowledge categories, 3 difficulty tiers, and 4 linguistic structures.
(1) A \emph{universal peak-depth constant} of $0.81\!\pm\!0.05$: correct answers
reliably peak near 83\% of model depth before final-layer collapse, with suppression
ratios $\rho$ reaching $201.5\times$.
(2) A sharp \emph{architectural dependency}: global and hybrid attention models
(GPT-2, Phi-2, Qwen, Mistral~7B) exhibit strong LLS while local and RoPE attention
models (GPT-Neo, Pythia, LLaMA~3~8B) show near-unity $\rho$.
The controlled 8B-scale comparison, Mistral~7B (hybrid, 25\% accuracy,
$\rho\!=\!47.6\times$) vs LLaMA~3~8B (RoPE, 70\% accuracy, $\rho\!=\!1.9\times$),
implicates attention architecture as a key driver of factual accuracy.
(3) Activation patching provides causal evidence in 8/8 cases
($\bar{\Delta p}\!=\!+0.774$).
(4) A zero-cost logit-blending intervention improves accuracy by up to $+45\%$ on
targeted benchmarks and \textbf{$+8.9$ percentage points on TruthfulQA (n=817)},
with zero effect on Type~2b knowledge-gap facts.
(5) Structural invariance: the peak-depth constant holds across 4 prompt structures
($0.77$--$0.83$).
(6) Multi-token facts show stronger suppression ($\rho\!=\!32.9\times$ vs
$11.6\times$), with non-Western facts showing equivalent profiles.
We release \textsc{HallScope} (\texttt{pip install hallscope}) and
\textsc{HallBench~v2}, a tiered benchmark distinguishing recoverable (Type~2a)
from knowledge-gap (Type~2b) hallucinations.
\end{abstract}

\noindent
\textbf{Code:} Code repository anonymised for review. \quad
\textbf{Library:} Analysis library anonymised for review. \quad
\textbf{Dataset:} Benchmark dataset anonymised for review.

\section{Introduction}

When GPT-2~XL is asked \textit{``The capital of France is''}, it outputs \textit{``a''}.
The standard explanation: the model never learned the correct answer.
Logit lens analysis tells a different story.
At block~41 of~48, ``Paris'' reaches 85.78\% probability.
By block~48, the final output layer, that probability has collapsed below 5\%.
The model encoded the correct answer. The final layer destroyed it.

We term this \textbf{Last-Layer Suppression} (LLS) and provide evidence
suggesting it is a consistent, causal, architecture-dependent mechanism
across 12 models spanning six families and two orders of magnitude in scale
(124M--8B parameters).

\paragraph{Contributions}
\begin{enumerate}[leftmargin=1.8em,topsep=2pt,itemsep=2pt,label=\arabic*.]
  \item \textbf{Universal peak-depth constant} $0.81\pm0.05$ across all 12 models
        (124M--8B), with $\rho$ reaching $201.5\times$.
  \item \textbf{Architectural dependency}: global/hybrid attention $\to$ strong LLS;
        local/RoPE attention $\to$ near-zero LLS.
        Mistral~7B vs LLaMA~3~8B (same scale, opposite profiles) provides a
        controlled 8B-scale comparison.
  \item \textbf{Causal evidence} via activation patching: 8/8 cases
        restored, $\bar{\Delta p}=+0.774$.
  \item \textbf{Zero-cost intervention}: $+45\%$ targeted, $+8.9$pp TruthfulQA
        (n=817), zero effect on Type~2b knowledge-gap facts.
  \item \textbf{Structural invariance}: depth constant across 4 prompt structures
        (680 total evaluation prompts).
  \item \textbf{Multi-token and cross-regional validity}: LLS stronger in
        multi-token facts; non-Western facts show equivalent suppression profiles.
  \item \textsc{HallScope} open-source library and \textsc{HallBench~v2} benchmark.
\end{enumerate}

\paragraph{Why this matters}
If hallucinations reflect suppressed rather than absent knowledge, then
training-time interventions targeting knowledge acquisition address the wrong
failure mode.
The Mistral/LLaMA comparison at 8B scale suggests that architectural choices
reducing LLS may be a primary driver of factual accuracy, a finding with direct
implications for model architecture decisions.

\section{Related Work}

\paragraph{Mechanistic interpretability}
The logit lens was introduced by Nostalgebraist~[1] and formalised as the tuned
lens by Belrose et al.~[2].
Elhage et al.~[3] established the residual stream framework that underlies our
analysis.
Meng et al.~[4] used causal tracing to localise factual associations in MLP
layers of GPT-J.
We extend these techniques to study suppression at the \emph{output} layer and
introduce systematic cross-architecture comparison at 7--8B scale.

\paragraph{Hallucination}
Ji et al.~[6] provide a comprehensive survey.
Lin et al.~[7] find that larger models are not necessarily more truthful on
TruthfulQA.
Manakul et al.~[8] detect hallucinations via consistency sampling.
Our approach targets the internal mechanism before generation and provides an
inference-time correction requiring no external data.

\paragraph{Knowledge and inference-time interventions}
Petroni et al.~[9] established that LLMs encode factual knowledge implicitly.
Geva et al.~[10] showed that MLP layers act as key-value memories.
Li et al.~[11] demonstrated representation engineering for truthfulness.
Our logit-blending intervention is simpler, no learned steering vectors, and
includes an explicit failure characterisation distinguishing Type~2a from
Type~2b cases.

\section{Background}

\paragraph{Logit lens}
A decoder-only transformer processes input tokens through $L$ layers, maintaining
a residual stream $\mathbf{h}^{(\ell)}\in\mathbb{R}^d$.
The logit lens applies the unembedding matrix $\mathbf{W}_U$ to intermediate states:
\begin{equation}
p_\ell(w) = \operatorname{softmax}\!\bigl(
  \mathrm{LN}(\mathbf{h}^{(\ell)})\cdot\mathbf{W}_U\bigr)[w]
\label{eq:lens}
\end{equation}
tracking the probability of any token $w$ at each layer $\ell=1,\ldots,L$.

\paragraph{Suppression ratio}
We define the suppression ratio $\rho$ as:
\begin{equation}
\rho = \frac{p_{\ell^*}(w^*)}{p_L(w^*)+\varepsilon},\quad
\ell^* = \operatorname*{arg\,max}_\ell\,p_\ell(w^*),\quad \varepsilon=10^{-10}
\label{eq:rho}
\end{equation}
where $w^*$ is the correct answer token and $p_L$ is the final-layer probability.
Values $\rho>2$ indicate meaningful suppression; $\rho>10$ is severe.

\paragraph{Logit blending}
\begin{equation}
\mathbf{logits}_{\mathrm{blend}}
  = \alpha\,\mathrm{LN}(\mathbf{h}^{(\ell^*)})\cdot\mathbf{W}_U
  + (1-\alpha)\,\mathrm{LN}(\mathbf{h}^{(L)})\cdot\mathbf{W}_U
\label{eq:blend}
\end{equation}
Setting $\alpha=0$ recovers standard generation.
The peak layer $\ell^*$ is identified per prompt via one additional forward pass.

\section{Methodology}

\paragraph{Models}
12 models across 6 families using TransformerLens~[5] and HuggingFace
Transformers.
\emph{Strong-suppression family} (global/hybrid attention):
GPT-2 Small (124M), Medium (345M), Large (774M), XL (1.5B);
Phi-2 (2.7B); Qwen~1.5-1.8B (1.8B); Mistral~7B.
\emph{Weak-suppression family} (local/RoPE attention):
GPT-Neo 125M, 1.3B, 2.7B; Pythia~2.8B; LLaMA~3~8B.
Mistral~7B and LLaMA~3~8B were loaded in 4-bit NF4 quantisation on
RTX~4060 (8~GB) and RTX~3080 (10~GB) respectively.

\paragraph{Hallucination taxonomy}
\textbf{Correct}: model outputs the correct answer.
\textbf{Type~2a} (LLS): correct token in top-10 at final layer but not argmax;
logit lens reveals an earlier probability peak followed by collapse.
\textbf{Type~2b} (Knowledge Gap): correct token absent from top-10 at all layers.

\paragraph{HallBench~v2}
150 facts across 10 categories (capitals, science, history, literature, modern,
geography, economy, culture) and 3 difficulty tiers (easy/medium/hard), including
30 non-Western facts (Africa, Asia, S.~America, Oceania).
Each fact is evaluated in 4 prompt structures:
forward completion, question form, reverse direction, and contextual embed.
This yields 600 structured prompts; 80 additional multi-token prompts bring the
total to 680 evaluation points.

\section{Experiments}

\subsection{The Peak-Depth Constant: Central Finding}

\begin{figure}[htbp]
  \centering
  \fig{figure1_layer_suppression}{Layer suppression curve}
  \caption{\textbf{Last-Layer Suppression in GPT-2~XL.}
  ``Paris'' peaks at 85.78\% probability at block~41
  (relative depth $\ell^*/L = 0.854$) before collapsing at the final
  layer block~48 ($\rho = 20.8\times$). The model encoded the correct
  answer; the final layer suppressed it.}
  \label{fig:suppression}
\end{figure}

\begin{figure}[htbp]
  \centering
  \fig{figure5_structure_variance}{Structure variance}
  \caption{\textbf{Peak-depth constant is invariant to prompt structure
  (GPT-2~XL, 30 facts $\times$ 4 structures = 120 prompts).}
  \textit{Left}: averaged layer-by-layer probability curves for all 4
  structures converge to the same suppression point near layer~40
  (dashed line = 0.83 depth constant).
  \textit{Right}: suppression ratio $\rho$ distributions by structure.
  The peak-depth constant across structures ($0.771$--$0.830$, range =
  0.059) is the most stable property of LLS. This structural invariance
  is the strongest evidence that LLS is a mechanistic property, not a
  prompt-format artefact.}
  \label{fig:structure}
\end{figure}

The central result of this paper is not the magnitude of $\rho$ in any
single case, which is subject to heavy-tailed outliers, but the
\emph{timing consistency} across all models and prompt forms.
Figure~\ref{fig:suppression} illustrates the phenomenon for a single case.
Figure~\ref{fig:structure} demonstrates its invariance across all 4 prompt
structures: forward, question, reverse direction, and contextual embed.
The peak-depth constant $0.81\pm0.05$ holds across all 12 models, all 4
structures, and both single- and multi-token facts.

\subsection{Cross-Architecture Comparison: 12 Models}

\begin{table}[htbp]
\caption{\textbf{Complete cross-model results across 12 models and 6
architecture families.}
Upper block: strong-suppression family (global/hybrid attention, shaded red).
Lower block: weak-suppression family (local/RoPE attention, shaded green).
The Mistral~7B / LLaMA~3~8B pair (same parameter count $\approx$8B,
same layer count = 32, highlighted in bold) provides a controlled
architectural comparison at scale.}
\label{tab:models}
\centering\small
\setlength{\tabcolsep}{5pt}
\begin{tabular}{llrrrrl}
\toprule
\textbf{Model} & \textbf{Family} & \textbf{Params}
  & \textbf{Acc} & \textbf{Type~2a} & \textbf{Avg $\rho$}
  & \textbf{Attention} \\
\midrule
\rowcolor{rowred}
GPT-2 Small    & GPT-2   & 124M & 10\% & 10\% & $2.4\times$  & Global \\
\rowcolor{rowred}
GPT-2 Medium   & GPT-2   & 345M & 20\% & 40\% & $15.6\times$ & Global \\
\rowcolor{rowred}
GPT-2 Large    & GPT-2   & 774M & 15\% & 15\% & $17.1\times$ & Global \\
\rowcolor{rowred}
GPT-2 XL       & GPT-2   & 1.5B & 35\% & 50\% & $20.8\times$ & Global \\
\rowcolor{rowred}
Phi-2          & Phi     & 2.7B & 70\% & 20\% & $10.8\times$ & Global \\
\rowcolor{rowred}
Qwen 1.5-1.8B  & Qwen    & 1.8B & 10\% & 70\% & $2.5\times$  & Global \\
\rowcolor{rowred}
\textbf{Mistral 7B} & \textbf{Mistral} & \textbf{7B}
  & \textbf{25\%} & \textbf{75\%} & $\mathbf{47.6\times}$ & GQA+SWA \\
\midrule
\rowcolor{rowgreen}
GPT-Neo 125M   & GPT-Neo & 125M &  5\% & 25\% & $1.0\times$  & Local \\
\rowcolor{rowgreen}
GPT-Neo 1.3B   & GPT-Neo & 1.3B & 40\% & 45\% & $1.0\times$  & Local \\
\rowcolor{rowgreen}
GPT-Neo 2.7B   & GPT-Neo & 2.7B & 15\% & 65\% & $1.0\times$  & Local \\
\rowcolor{rowgreen}
Pythia 2.8B    & Pythia  & 2.8B & 40\% & 50\% & $1.1\times$  & Rotary \\
\rowcolor{rowgreen}
\textbf{LLaMA 3 8B} & \textbf{LLaMA~3} & \textbf{8B}
  & \textbf{70\%} & \textbf{30\%} & $\mathbf{1.9\times}$
  & RoPE+GQA \\
\bottomrule
\end{tabular}
\end{table}

\begin{figure}[htbp]
  \centering
  \fig{figure3_architecture_families}{Architecture families}
  \caption{\textbf{Two hallucination families across 12 models.}
  Global/hybrid attention models (red region) cluster at high $\rho$.
  Local/RoPE models (green region) cluster near $\rho\approx1$.
  Suppression ratio predicts intervention effectiveness ($r=0.91$,
  $p<0.001$). The Mistral~7B / LLaMA~3~8B pair (circled) provides a
  controlled comparison at equal scale, isolating the architectural effect.}
  \label{fig:families}
\end{figure}

\begin{figure}[htbp]
  \centering
  \fig{figure2_scaling_law}{Scaling law}
  \caption{\textbf{Scaling laws within the GPT-2 family (124M--1.5B).}
  \textit{Left}: suppression ratio $\rho$ increases monotonically with
  scale ($2.4\times$ to $20.8\times$).
  \textit{Centre}: Type~2a rate increases (10\% to 50\%).
  \textit{Right}: peak relative depth converges to $\approx0.83$ across
  all four sizes.}
  \label{fig:scaling}
\end{figure}

Table~\ref{tab:models} and Figure~\ref{fig:families} reveal a clean
architectural dichotomy.
The key controlled comparison is \textbf{Mistral~7B vs LLaMA~3~8B}:
both have approximately 8B parameters and 32 layers, trained on modern
large-scale data.
Mistral~7B uses Sliding Window Attention with global-capable layers:
Type~2a rate \textbf{75\%} (highest of any model), avg $\rho=47.6\times$,
accuracy \textbf{25\%}.
LLaMA~3~8B uses pure RoPE positional encoding:
Type~2a rate 30\%, avg $\rho=1.9\times$, accuracy \textbf{70\%}
(highest of any model).
At equal scale, the model with lower suppression has nearly $3\times$
better factual accuracy, strongly implicating attention architecture.

A further pattern: within the strong-suppression family, suppression
severity \emph{increases} with model capability.
GPT-2~XL ($\rho=20.8\times$) $\to$ Mistral~7B ($\rho=47.6\times$).
More capable models encode facts more strongly in intermediate layers
while suppressing them more aggressively at the output, a counterintuitive
finding with implications for scaling-based hallucination reduction strategies.

\subsection{Causal Evidence via Activation Patching}

\begin{table}[htbp]
\caption{\textbf{Self-patching on GPT-2~XL.} Injecting peak-layer
activations into the final layer restores correct predictions in all
8 tested cases. The Australia case (already correct) shows
$\Delta p=+0.829$, demonstrating that even correct predictions are
suppressed but survive because their signal exceeds the suppression
magnitude.}
\label{tab:patching}
\centering\small
\begin{tabular}{lccr}
\toprule
\textbf{Prompt} & \textbf{Baseline} & \textbf{Patched} & $\Delta p$ \\
\midrule
The capital of France is     & \textit{a}   & Paris      & $+0.812$ \\
The capital of Russia is     & \textit{the} & Moscow     & $+0.901$ \\
The capital of Germany is    & \textit{a}   & Berlin     & $+0.716$ \\
The capital of China is      & \textit{a}   & Beijing    & $+0.725$ \\
The capital of India is      & \textit{a}   & Delhi      & $+0.801$ \\
Evolution proposed by        & Charles      & Darwin     & $+0.903$ \\
Einstein discovered          & \textit{the} & relativity & $+0.505$ \\
The capital of Australia is  & Canberra     & Canberra   & $+0.829$ \\
\midrule
\textbf{Mean}               & ---          & ---        & $\mathbf{+0.774}$ \\
\bottomrule
\end{tabular}
\end{table}

Self-patching injects peak-layer ($\ell^*$) activations from the \emph{same}
prompt into the final layer, bypassing suppression.
Table~\ref{tab:patching} shows restoration in 8/8 cases with mean
$\bar{\Delta p}=+0.774$.
This constitutes strong causal evidence suggesting that the final
transformer layers actively suppress factual representations that are
present and strong in intermediate layers.

\subsection{Logit Blending Intervention}

\begin{table}[htbp]
\caption{\textbf{Logit blending results.} Strong-suppression models improve
substantially; weak-suppression models show zero effect.
The TruthfulQA / TriviaQA contrast confirms the Type~2a/2b taxonomy.}
\label{tab:intervention}
\centering\small
\begin{tabular}{llrrrr}
\toprule
\textbf{Model / Benchmark} & \textbf{Type}
  & \textbf{Baseline} & $\alpha{=}0.3$ & $\alpha{=}0.5$ & \textbf{Best} \\
\midrule
GPT-2 XL (targeted)    & Strong & 40\%   & $+35\%$ & $+45\%$ & 85\%   \\
Qwen 1.5-1.8B          & Strong & 10\%   & $+35\%$ & $+40\%$ & 50\%   \\
GPT-Neo 2.7B           & Weak   & 15\%   & $+0\%$  & $+0\%$  & 15\%   \\
LLaMA 3 8B             & Weak   & 70\%   & $+0\%$  & $+0\%$  & 70\%   \\
\midrule
TruthfulQA (n=817)     & Mixed  & 20.8\% & $+6.4\%$& $+8.9\%$& 29.7\% \\
TriviaQA (n=1000)      & Type~2b&  3.5\% & $+0.0\%$& $+0.1\%$&  3.6\% \\
\bottomrule
\end{tabular}
\end{table}

\paragraph{TruthfulQA vs TriviaQA: taxonomy confirmed}
The contrast between these benchmarks is the strongest population-level
evidence for the Type~2a/2b distinction.
TruthfulQA questions probe facts that models encode but suppress (Type~2a
dominated): $+8.9$ percentage points on 817 questions.
TriviaQA's 3.5\% baseline indicates knowledge-gap dominance (Type~2b):
$+0.1$pp, exactly as predicted, the intervention cannot manufacture
absent knowledge.
This bidirectional confirmation distinguishes a mechanistic claim from a
benchmark-tuning result.

\paragraph{Failure profile}
On 15 Type~2b prompts: 0/15 correct at all $\alpha$ values.
On 10 non-factual prompts: 3/10 predictions change at $\alpha=0.5$,
all semantically coherent.
Alpha sweep ($\alpha\in[0,1]$, step 0.1): accuracy rises monotonically
from 50\% to 90\% at $\alpha\geq0.5$, no degradation cliff above optimal.

\subsection{Multi-Token and Cross-Regional Validity}

\begin{figure}[htbp]
  \centering
  \fig{figure6_multi_token}{Multi-token analysis}
  \caption{\textbf{Multi-token and cross-regional analysis (GPT-2~XL,
  42 facts).}
  \textit{Left}: multi-token facts show \emph{higher} average $\rho$
  ($32.9\times$) than single-token ($11.6\times$), refuting the
  hypothesis that LLS is a single-token artefact.
  \textit{Centre}: hallucination type distribution by token count.
  \textit{Right}: $\rho$ by geographic region, non-Western facts
  (Africa, Asia, S.~America) show suppression profiles comparable to
  Western facts.}
  \label{fig:multitoken}
\end{figure}

\begin{table}[htbp]
\caption{\textbf{Single-token vs multi-token suppression} (GPT-2~XL,
Experiment~12). Multi-token suppression is stronger, not weaker.}
\label{tab:multitoken}
\centering\small
\begin{tabular}{lrrrrrr}
\toprule
\textbf{Type} & \textbf{n} & \textbf{Acc}
  & \textbf{Type~2a} & \textbf{Avg $\rho$}
  & \textbf{Type~2a $\rho$} & \textbf{Depth} \\
\midrule
Single-token & 22 & 50\% & 36\% & $11.6\times$ & $14.9\times$ & 0.836 \\
Multi-token  & 20 & 25\% & 45\% & $32.9\times$ & $20.6\times$ & 0.912 \\
\bottomrule
\end{tabular}
\end{table}

Multi-token facts show \emph{higher} Type~2a rates (45\% vs 36\%) and
higher average $\rho$ ($32.9\times$ vs $11.6\times$) than single-token
facts.
The depth constant shifts slightly deeper (0.912 vs 0.836) but remains
within $0.81\pm0.05$.
Extreme cases: yen ($\rho=386.4\times$), euro ($86.1\times$),
Rowling ($82.3\times$).
LLS is stronger, not weaker, for multi-token factual associations,
directly refuting the single-token-artefact critique.

\subsection{HallBench~v2 and Closed Model Probing}

\begin{figure}[htbp]
  \centering
  \fig{figure4_hallbench_v2}{HallBench v2}
  \caption{\textbf{HallBench~v2 tiered results.}
  Tier~1 (high suppression facts: capitals, science) shows the largest
  intervention gains (GPT-2~XL: $36\%\to80\%$).
  Tier~3 (knowledge gap) shows \emph{zero} improvement across all models
  at all $\alpha$ values, the intervention never manufactures absent
  knowledge.
  This tier separation is the strongest per-fact evidence that logit
  blending targets the suppression mechanism specifically.}
  \label{fig:hallbench}
\end{figure}

\paragraph{Closed model probing}
Behavioural probing of Gemini~2.5 Flash (no internal access):
\textit{``The Berlin Wall fell in''} elicits \textit{``19''} while
direct question mode elicits \textit{``1989''}.
Only 1/4 phrasings yield the complete year; three variants produce
truncated outputs.
This systematic truncation is consistent with LLS cutting off the
complete factual token at generation time, though mechanistic equivalence
cannot be confirmed without internal access.

\paragraph{Cross-lingual finding}
Activation patching of Qwen~1.5-1.8B on
\textit{``The capital of Russia is''} yields
\begin{CJK}{UTF8}{gbsn}莫斯科\end{CJK} (Moscow in Chinese)
rather than the English token.
This reveals cross-lingual factual encoding normally suppressed by both
LLS and a language-switch mechanism at the output layer.

\section{The \textsc{HallScope} Library}

\texttt{pip install hallscope} provides logit lens analysis, Type~2a/2b
classification, suppression ratio computation, activation patching, logit
blending, and HallBench~v2 integration.

\begin{verbatim}
from hallscope import HallScope
hs = HallScope("gpt2-xl")
report = hs.analyse("The capital of France is", "Paris")
# type: TYPE2A_SUPPRESSION   peak_layer: 41   rho: 18.6x
corrected = hs.correct("The capital of France is")
# "Paris"
\end{verbatim}

\section{Discussion}

\paragraph{Why does LLS occur?}
The Mistral~7B / LLaMA~3~8B comparison at equal scale provides the
clearest evidence: attention architecture determines suppression profile.
Sliding Window Attention retains global-capable layers that may learn to
prioritise distributional patterns over factual content at the output.
Pure RoPE positional encoding limits this global override.
This hypothesis requires circuit-level ablation for confirmation;
confounds include training data differences between model families.

\paragraph{Suppression increases with capability}
Within the strong-suppression family, larger and more capable models
suppress more aggressively: GPT-2~XL ($20.8\times$) $\to$ Mistral~7B
($47.6\times$).
If this trend continues at frontier scale, scaling-based approaches to
hallucination reduction may inadvertently worsen the retrieval problem.

\paragraph{Implications for knowledge editing}
Methods such as ROME~[4] and MEMIT address factual storage in MLP layers.
LLS suggests that even correctly stored facts can be suppressed at the
output layer, knowledge editing may address the storage problem while
leaving the retrieval problem intact.
We leave direct experimental comparison with ROME and MEMIT for future
work; they target the storage mechanism while our intervention targets
inference-time retrieval, and a fair comparison would require a shared
evaluation protocol that distinguishes Type~2a from Type~2b failures.

\paragraph{Type~2a vs Type~2b diagnostics}
Standard benchmarks cannot distinguish between suppression and absent
knowledge.
\textsc{HallBench~v2} is designed to make this distinction measurable.
We recommend that future hallucination research report Type~2a and Type~2b
rates separately to enable targeted mitigation strategies.

\section{Limitations}

\textbf{Scale ceiling.} Our largest models are 7--8B parameters.
Whether LLS persists, intensifies, or diminishes at frontier scale
($>70$B) is unknown.
The trend within the strong-suppression family (suppression worsens with
scale) is concerning but requires validation.

\textbf{Logit lens assumptions.} Applying the final-layer unembedding to
intermediate states assumes representations are already in unembedding
space.
The tuned lens~[2] addresses this with per-layer learned probes.
Our depth-constant finding is robust since it measures timing rather
than absolute probability magnitude.

\textbf{Sample size per model.} Individual model analyses use 20 prompts.
Results are indicative; large-scale validation would strengthen claims.
The structural invariance experiment (120 prompts) and benchmark evaluations
(TruthfulQA n=817, TriviaQA n=1000) provide larger-scale validation.

\textbf{Closed model evidence.}
Gemini~2.5 Flash results are behavioural only; mechanistic equivalence
with open-model LLS cannot be confirmed.

\textbf{Multi-token intervention.}
Logit blending addresses single-step next-token prediction.
Multi-token answers require autoregressive generation that a single-step
intervention does not correct.

\textbf{Architectural confounds.}
Training data, tokenizer, and procedure differ across architecture families.
Isolating attention mechanism as the sole cause requires controlled
ablation beyond the scope of this paper.

\section{Conclusion}

Our findings suggest that Last-Layer Suppression is a consistent, causal,
and architecture-dependent mechanism of factual hallucination across 12
models spanning 6 families and two orders of magnitude (124M--8B).
The peak-depth constant $0.81\pm0.05$ holds across all models, all prompt
structures, and both token-count classes.
Activation patching provides causal evidence (8/8 cases, $\bar{\Delta p}=+0.774$).
Logit blending recovers suppressed knowledge without manufacturing absent
knowledge ($+8.9$pp on TruthfulQA, $+0.1$pp on Type~2b-dominated TriviaQA).
The Mistral~7B / LLaMA~3~8B comparison at equal scale, the best open model
(70\% accuracy, $\rho=1.9\times$) vs the most heavily suppressed 7B model
(25\% accuracy, $\rho=47.6\times$), suggests that architectural choices
affecting LLS may be a key driver of factual accuracy in large language models.

These findings reframe the hallucination problem: for a substantial fraction
of errors, the model is not ignorant, it is suppressing knowledge it has
already encoded.
Addressing hallucination may require as much attention to the retrieval
process at the output layer as to the storage process in intermediate layers.

\section*{Acknowledgements}

Withheld for double-blind review.

\section*{References}

{\small
\begin{enumerate}[leftmargin=2em,itemsep=2pt,label={[\arabic*]}]

\item Nostalgebraist (2020).
  Interpreting GPT: the logit lens.
  \textit{LessWrong}.
  \url{https://www.lesswrong.com/posts/AcKRB8wDpdaN6v6ru}

\item N.\ Belrose, Z.\ Furman, L.\ Smith, D.\ Halawi, I.\ Ostrovsky,
  L.\ McKinney, S.\ Biderman, \& J.\ Steinhardt (2023).
  Eliciting latent predictions from transformers with the tuned lens.
  \textit{arXiv:2303.08112}.

\item N.\ Elhage, N.\ Nanda, C.\ Olsson, T.\ Henighan, N.\ Joseph,
  B.\ Mann et al.\ (2021).
  A mathematical framework for transformer circuits.
  \textit{Transformer Circuits Thread}.

\item K.\ Meng, D.\ Bau, A.\ Andonian, \& Y.\ Belinkov (2022).
  Locating and editing factual associations in GPT.
  \textit{Advances in Neural Information Processing Systems (NeurIPS)}, 35.

\item N.\ Nanda (2022).
  TransformerLens: A library for mechanistic interpretability.
  \textit{GitHub}.
  \url{https://github.com/neelnanda-io/TransformerLens}

\item Z.\ Ji, N.\ Lee, R.\ Frieske, T.\ Yu, D.\ Su, Y.\ Xu et al.\ (2023).
  Survey of hallucination in natural language generation.
  \textit{ACM Computing Surveys}, 55(12), 1--38.

\item S.\ Lin, J.\ Hilton, \& O.\ Evans (2022).
  TruthfulQA: Measuring how models mimic human falsehoods.
  \textit{Proceedings of ACL 2022}.

\item P.\ Manakul, A.\ Liusie, \& M.\ J.\ F.\ Gales (2023).
  SelfCheckGPT: Zero-resource black-box hallucination detection for
  generative large language models.
  \textit{Proceedings of EMNLP 2023}.

\item F.\ Petroni, T.\ Rockt\"{a}schel, P.\ Lewis, A.\ Bakhtin,
  Y.\ Wu, A.\ Miller, \& S.\ Riedel (2019).
  Language models as knowledge bases?
  \textit{Proceedings of EMNLP 2019}.

\item M.\ Geva, R.\ Schuster, J.\ Berant, \& O.\ Levy (2021).
  Transformer feed-forward layers are key-value memories.
  \textit{Proceedings of EMNLP 2021}.

\item K.\ Li, O.\ Patel, F.\ Vi\'{e}gas, H.\ Pfister, \& M.\ Wattenberg
  (2023).
  Inference-time intervention: Eliciting truthful answers from a language
  model.
  \textit{Advances in Neural Information Processing Systems (NeurIPS)}, 36.

\item S.\ Kadavath, T.\ Conerly, A.\ Askell, T.\ Henighan, D.\ Drain,
  E.\ Perez et al.\ (2022).
  Language models (mostly) know what they know.
  \textit{arXiv:2207.05221}.

\item A.\ Zou, L.\ Phan, S.\ Chen, J.\ Campbell, P.\ Guo, R.\ Ren et al.\
  (2023).
  Representation engineering: A top-down approach to AI transparency.
  \textit{arXiv:2310.01405}.

\item E.\ Hernandez, A.\ Ghandeharioun, L.\ Shen, Z.\ Li, R.\ Ghosh,
  R.\ Kim et al.\ (2023).
  Linearity of relation decoding in transformer language models.
  \textit{Proceedings of ICLR 2024}.

\end{enumerate}
}

\newpage
\section*{NeurIPS Paper Checklist}

\noindent\textit{The checklist is included as required by the NeurIPS 2026
Call for Papers. It does not count toward the page limit.}

\vspace{6pt}

{\small
\begin{enumerate}[leftmargin=*,label=\arabic*.,itemsep=8pt]

\item \textbf{Claims.}
\textit{Do the main claims made in the abstract and introduction accurately
reflect the paper's contributions and scope?}

\textbf{Yes.}
All 7 contributions are enumerated explicitly in the introduction.
Claims use calibrated language throughout: ``consistent'' not ``universal'';
``strongly implicates'' not ``proves'' for the architectural hypothesis.
The Limitations section (Section~8) explicitly bounds scope to $\leq$8B
parameters, 20 prompts per individual model analysis, and notes architectural
confounds.

\item \textbf{Limitations.}
\textit{Did you discuss the limitations of your work?}

\textbf{Yes.}
Section~8 provides six explicit limitations:
(1) scale ceiling, behaviour at $>70$B parameters is unknown;
(2) logit lens intermediate-state assumption;
(3) sample size per model (20 prompts for individual analysis);
(4) closed model evidence is behavioural only;
(5) multi-token intervention scope limited to single-step prediction; and
(6) architectural confounds including training data differences.

\item \textbf{Theory, Assumptions and Proofs.}
\textit{Did you state all assumptions and include complete proofs?}

\textbf{N/A.}
The paper makes no formal theoretical claims requiring proof.
Equations~1--3 state all mathematical definitions with explicit assumptions
(the $\varepsilon$ regularisation in Eq.~2; the per-prompt forward pass
for $\ell^*$ in Eq.~3).
No theorems or lemmas are claimed.

\item \textbf{Experimental Result Reproducibility.}
\textit{Did you include code, data, and instructions?}

\textbf{Yes.}
All 16 experiment scripts, the \textsc{HallScope} analysis library,
and the \textsc{HallBench~v2} dataset will be released upon acceptance.
An anonymised version is included in the supplementary material.
All base models used are publicly available via HuggingFace.
Each experiment script is self-contained with exact commands and
environment specifications documented in the supplementary material.

\item \textbf{Open Access to Data and Code.}
\textit{Did you include the code, data, and instructions needed to reproduce
the main experimental results?}

\textbf{Yes.}
Code: \texttt{[anonymised]}\\
Library: \texttt{[anonymised]}\\
Dataset: \texttt{[anonymised]}\\
All experiment scripts, the \textsc{HallScope} library, and the
\textsc{HallBench~v2} dataset will be released upon acceptance.
An anonymised version of the code and data is included in the
supplementary material submitted with this paper.

\item \textbf{Experimental Setting/Details.}
\textit{Did you specify all training and evaluation details?}

\textbf{Yes.}
Section~4 specifies all 12 models with parameter counts, architecture
families, and quantisation settings (4-bit NF4 for 7B+ models via
BitsAndBytes).
Hardware is stated: RTX~4060 (8~GB VRAM) and RTX~3080 (10~GB VRAM).
Classification criteria are defined in Section~4: Type~2a requires
correct token rank $\leq10$ at the final layer; Type~2b requires rank
$>10$.
The intervention alpha sweep range $[0,1]$ with step 0.1 is documented
in Section~5.4.

\item \textbf{Experiment Statistical Significance.}
\textit{Did you report error bars or statistical significance?}

\textbf{No.}
Error bars are not reported for individual model analyses (20 prompts per
model is insufficient for formal significance testing).
We acknowledge this in Section~8.
However, three factors provide qualitative strength:
(a) the structural invariance experiment uses 120 prompts across 4
structures with consistent results;
(b) TruthfulQA (n=817) and TriviaQA (n=1000) provide large-sample
benchmark validation;
(c) the Mistral~7B / LLaMA~3~8B controlled comparison isolates the
architectural effect at equal scale.

\item \textbf{Experiments Compute Resources.}
\textit{Did you provide information on compute resources?}

\textbf{Yes.}
All experiments ran on consumer hardware: RTX~4060 Laptop GPU (8~GB VRAM)
for GPT-2 family, Phi-2, Qwen, and Mistral~7B (4-bit quantisation);
RTX~3080 (10~GB VRAM) for LLaMA~3~8B (4-bit quantisation).
No cloud compute was used.
TransformerLens experiments complete in under 2 hours per model;
7B+ models require approximately 30 minutes per model after download.
Total estimated GPU-hours: approximately 20 hours across all 16 experiments.

\item \textbf{Code of Ethics.}
\textit{Have you read the NeurIPS Code of Ethics and ensured that your
research conforms to it?}

\textbf{Yes.}
The research analyses existing publicly available open-source language
models.
No human subjects, personal data, or harmful applications are involved.
The primary aim, reducing factual hallucinations in language models, is
a beneficial goal for AI safety and reliability.

\item \textbf{Broader Impacts.}
\textit{Did you discuss potential negative societal impacts?}

\textbf{Yes.}
The logit-blending intervention requires white-box access to model internals
(residual stream activations), which limits practical misuse potential.
The primary societal impact is positive: reducing hallucinations in AI
systems deployed in high-stakes domains such as healthcare, legal, and
educational applications.
A theoretical dual-use risk is selective suppression or surfacing of
specific knowledge; this is noted as a precaution but requires internal
model access and is unlikely to represent a significant new threat vector.

\item \textbf{Safeguards.}
\textit{Do you have safeguards for responsible model release?}

\textbf{N/A.}
No new pretrained language models are released.
\textsc{HallScope} is an analysis library containing no harmful content.
\textsc{HallBench~v2} is a factual benchmark dataset containing no harmful
content.
All analysed models are already publicly available through their respective
model pages on HuggingFace.

\item \textbf{Licenses.}
\textit{Did you cite creators and respect licenses?}

\textbf{Yes.}
All models and tools are cited in References.
Licenses: GPT-2 (MIT), Phi-2 (MIT), Qwen~1.5 (Tongyi Qianwen License),
Mistral~7B (Apache~2.0), GPT-Neo and Pythia (Apache~2.0), LLaMA~3~8B
(Meta LLaMA~3 Community License), TransformerLens (MIT License).
TruthfulQA and TriviaQA are used for evaluation only under their respective
research licences.
Released assets: \textsc{HallScope} (MIT License),
\textsc{HallBench~v2} (CC-BY~4.0).

\item \textbf{Assets.}
\textit{Did you document new assets?}

\textbf{Yes.}
\textsc{HallScope} (anonymised, MIT License) provides
logit lens analysis, suppression ratio computation, Type~2a/2b classification,
activation patching, logit blending, and HallBench~v2 integration.
Full API documentation is available in the GitHub repository.
\textsc{HallBench~v2} (HuggingFace: \texttt{[anonymised]},
CC-BY~4.0) contains 50 curated examples with tier labels, first-token
answers, suppression metadata, and source attribution.
No personal data is included in either release.

\item \textbf{Crowdsourcing and Research with Human Subjects.}
\textit{Did you include full details of human subject research?}

\textbf{N/A.}
No crowdsourcing or human subjects research was conducted.
All experiments are fully computational analyses of existing open-source
language models on existing benchmark datasets.

\item \textbf{IRB Approvals.}
\textit{Did you obtain IRB approval?}

\textbf{N/A.}
No human subjects research was conducted. IRB approval is not applicable.

\item \textbf{Declaration of LLM Usage.}
\textit{Does the paper describe LLM usage as an important component of
core methods?}

\textbf{N/A.}
Large language models are the \emph{subject of study} in this paper,
not components of the methodology.
The analysis pipeline uses TransformerLens and HuggingFace Transformers
for mechanistic analysis; no LLM is used to generate scientific content,
write experiment code, or make analytical decisions in this work.

\end{enumerate}
}

\end{document}